% section: 手法の見極め

\section{手法の見極め}

漸化式を見たら,いきなり式変形を始める前に,次の順で観察する.

\subsection{第1段階: 定義域と添字を確認する}

分母が $0$ にならないか,対数の中身が正か,平方根の中身が非負かを
確認する.また,式が $a_0$ から始まるのか $a_1$ から始まるのかを
固定する.この確認を省くと,形式的には正しい変形でも数列の途中で
意味を失うことがある.

\subsection{第2段階: 変化量を計算する}

次の三つを調べる.
\[
  a_{n+1}-a_n,
  \qquad
  \frac{a_{n+1}}{a_n},
  \qquad
  a_{n+1}=F(n,a_n).
\]
差が $n$ の式だけなら階差型,比が $n$ の式だけなら階比型である.
差や比に整理しても $a_n$ が残るなら,その残り方に注目する.

\subsection{第3段階: 次の候補を試す}

\begin{description}
  \item[定数が足される] 不動点を求め,$a_n-\alpha$ を考える.
  \item[$f(n)$ が残る] 補助数列を設計し,右辺の非斉次項を消す.
  \item[積やべきが現れる] 対数を取る.ただし,和が残るなら対数だけでは
    解決しない.
  \item[分数が現れる] 逆数,不動点,一次分数変換のいずれが自然かを比べる.
  \item[過去2項以上が現れる] 特性方程式,行列,または差分演算子を使う.
  \item[多項式の反復が現れる] 三角関数や対称式など,恒等式で形が閉じる
    表現を探す.
  \item[畳み込みが現れる] 母関数に移して,和を積へ翻訳する.
  \item[どの形にも入らない] 一般項を一度保留し,不動点,単調性,有界性,
    周期性,近似誤差を調べる.
\end{description}

\subsection{第4段階: 候補を検証する}

候補を思いついたら,初項と漸化式の二つを確認する.初項だけに一致する
式は,数列の一般項とは限らない.漸化式だけを形式的に満たす式も,
初期条件が違えば別の数列である.

帰納法で確認する場合は,次の三行を明確に書く.
\begin{enumerate}
  \item $n=1$ または $n=0$ で候補が初項に一致する.
  \item $a_n$ が候補の形だと仮定する.
  \item 漸化式へ代入すると $a_{n+1}$ が同じ候補の $n+1$ の形になる.
\end{enumerate}

この確認まで終わって,初めて「解けた」と言える.公式との類似を指摘
するだけでは,方針であって解答ではない.
